\documentclass[11pt]{article}
\usepackage[T1]{fontenc}
\usepackage{lmodern}
\usepackage{amsmath,amssymb,amsthm}
\usepackage[margin=1in]{geometry}
\usepackage[colorlinks=true,linkcolor=blue,citecolor=blue,urlcolor=blue]{hyperref}

\newtheorem{theorem}{Theorem}
\newtheorem{remark}{Remark}

\newcommand{\spn}{\operatorname{span}}
\newcommand{\clspn}{\overline{\operatorname{span}}}
\newcommand{\Krho}{K_\varrho}
\newcommand{\Frho}{\mathcal F_\varrho}
\newcommand{\Xrho}{X_\varrho}

\title{Coordinate Subspaces of the HRST Norming-M-Basis Space}
\author{Run \texttt{fa\_banach\_001}, agent \texttt{agent\_lane\_12}}
\date{2026-06-15}

\begin{document}
\maketitle

\section*{Source Problem}

Hajek and Russo ask whether admitting a norming Markushevich basis is
hereditary to subspaces \cite[Problem 6.1]{HR}. The problem appears on PDF
page 24 of arXiv:2402.18442 and in the parsed source at
\texttt{manuscript.tex:685--691}.

The same source discusses the Hajek--Russo--Somaglia--Todorcevic construction:
there is an Asplund space $\Xrho$ with a $1$-norming M-basis which is not WCG
\cite[Theorem A, recalled as Theorem 4.4]{HR}. This makes $\Xrho$ a natural
place to search for a possible bad subspace in Problem 6.1.

\section*{Result}

The coordinate-generated subspaces of $\Xrho$ are not counterexamples.

\begin{theorem}\label{thm:coordinate-subspaces}
Let $\Xrho=\clspn\{f_\alpha:\alpha<\omega_1\}\subset C(\Krho)$ be the HRST
space as recalled in \cite[Theorem 4.4]{HR}. For every
$A\subseteq\omega_1$, put
\[
    X_A=\clspn\{f_\alpha:\alpha\in A\}\subseteq \Xrho .
\]
Then
\[
    \{f_\alpha;\ \mu_\alpha|_{X_A}\}_{\alpha\in A}
\]
is a $1$-norming M-basis of $X_A$.
\end{theorem}

\section*{Proof Intuition}

The HRST construction is unusually rigid on its coordinate subspaces. The
compact $\Krho$ is the closure of a family $\Frho$ of finite subsets of
$\omega_1$. On $\Xrho$, evaluation at any finite set
$F\in\Frho$ is exactly the finite sum of the coordinate measures indexed by
$F$. If we keep only coordinates in $A$, the same evaluation is simply the
finite sum over $F\cap A$. Since the finite sets $\Frho$ are dense in
$\Krho$, those evaluations still norm the restricted space.

\section*{Proof}

\begin{proof}
Recall the notation from the source paper. The compact space $\Krho$ is the
pointwise closure of
\[
    \Frho=\{F_n(\alpha):n<\omega,\ \alpha<\omega_1\}
    \subseteq [\omega_1]^{<\omega}.
\]
For $\alpha<\omega_1$, the coordinate function $f_\alpha\in C(\Krho)$ is
defined by
\[
    f_\alpha(B)=
    \begin{cases}
    1, & \alpha\in B,\\
    0, & \alpha\notin B,
    \end{cases}
    \qquad B\in\Krho,
\]
and $\mu_\alpha=\delta_{\{\alpha\}}$ is the point mass at the singleton
$\{\alpha\}\in\Krho$. The system
$\{f_\alpha;\mu_\alpha|_{\Xrho}\}_{\alpha<\omega_1}$ is the canonical
$1$-norming M-basis of $\Xrho$ in the HRST theorem.

Fix $A\subseteq\omega_1$ and define $X_A$ as in the statement. The displayed
system is biorthogonal because, for $\alpha,\beta\in A$,
\[
    \mu_\alpha(f_\beta)=\delta_{\{\alpha\}}(f_\beta)=f_\beta(\{\alpha\})
    =\delta_{\alpha\beta}.
\]
It is fundamental by the definition of $X_A$.

It remains to prove the $1$-norming property. Let $F\in\Frho$. Since $F$ is
finite, the evaluation functional $\delta_F|_{X_A}$ satisfies
\[
    \delta_F|_{X_A}
      =
    \sum_{\alpha\in F\cap A}\mu_\alpha|_{X_A}.
\]
Indeed, both sides agree on every generator $f_\beta$ with $\beta\in A$:
they both give $1$ if $\beta\in F$ and $0$ otherwise. Hence the equality
holds on the dense linear span of the generators and therefore on $X_A$.
Moreover, $\|\delta_F|_{X_A}\|\le 1$.

For $x\in X_A$, continuity on the compact space $\Krho$ and density of
$\Frho$ in $\Krho$ give
\[
    \|x\|_{C(\Krho)}
      =
    \sup_{B\in\Krho}|x(B)|
      =
    \sup_{F\in\Frho}|\delta_F(x)|.
\]
By the preceding identity, each $\delta_F|_{X_A}$ belongs to
$\spn\{\mu_\alpha|_{X_A}:\alpha\in A\}$ and has norm at most one. Therefore
\[
    \|x\|
    \le
    \sup\Big\{
        |\nu(x)|:\nu\in\spn\{\mu_\alpha|_{X_A}:\alpha\in A\},
        \ \|\nu\|\le 1
    \Big\}.
\]
The reverse inequality is immediate from the condition $\|\nu\|\le1$.
Thus the span of the restricted coordinate measures is $1$-norming for
$X_A$. In particular it is total, and the biorthogonal fundamental system is
a $1$-norming M-basis.
\end{proof}

\section*{Limitations and Novelty Check}

This theorem does not settle the full heredity problem. It only treats
subspaces of $\Xrho$ that are generated by subsets of the canonical HRST
coordinate family. A bad subspace of $\Xrho$, if one exists, must therefore
mix the coordinates in an essentially non-coordinate way.

The result was obtained by re-reading the HRST norming argument as presented
in the source survey. It appears to be an immediate but useful consequence of
the finite-set evaluation formula in that construction. No claim is made that
this observation is absent from the original HRST paper; its value here is as
a recorded search filter for Problem 6.1.

\section*{Human Review Recommendation}

Review as a positive coordinate-subspace subcase for Problem 6.1 and as a
counterexample-search obstruction inside the HRST example. The proof is
short; the main points are the equality
$\delta_F|_{X_A}=\sum_{\alpha\in F\cap A}\mu_\alpha|_{X_A}$ and the use of
density of $\Frho$ in $\Krho$.

\begin{thebibliography}{9}

\bibitem{HR}
P. Hajek and T. Russo,
\emph{Norming Markushevich bases: recent results and open problems},
arXiv:2402.18442.

\end{thebibliography}

\end{document}
